\documentclass[12pt]{ctexart}

\usepackage[a4paper,margin=2.5cm]{geometry}
\usepackage{amsmath,amssymb}
\usepackage{enumitem}

\setlist[enumerate]{label=(\arabic*),leftmargin=2em}

\title{截图内容提取}
\author{}
\date{}

\begin{document}

\maketitle

\section*{圆锥曲线——点乘双根法}

\subsection*{适用类型}

类似
\[
x_1x_2,\qquad y_1y_2,\qquad (x_1+t)(x_2+t),\qquad (y_1+t)(y_2+t),
\]
或
\[
\overrightarrow{MA}\cdot\overrightarrow{MB},\qquad
|\overrightarrow{MA}|\cdot|\overrightarrow{MB}|
\]
的结构。

其中 $x_1,x_2,y_1,y_2$ 是直线与曲线的两个交点的横纵坐标，
$A,B$ 是直线与曲线的两个交点。该方法也适用于可以转化为上述结构的问题。

\subsection*{理论基础}

二次函数的双根式。若一元二次方程
\[
ax^2+bx+c=0\qquad (a\ne 0)
\]
的两根为 $x_1,x_2$，则
\[
ax^2+bx+c=a(x-x_1)(x-x_2).
\]

\subsection*{具体步骤}

\[
\text{化双根式}\longrightarrow \text{赋值}\longrightarrow \text{变形代入}.
\]

\section*{例题 1}

椭圆
\[
C:\frac{x^2}{4}+\frac{y^2}{3}=1
\]
左顶点为 $A$，过左焦点 $F$ 的直线交椭圆于 $P,Q$ 两点。

求证：直线 $AP$ 和 $AQ$ 的斜率乘积是定值。

\section*{例题 2}

（2021 新高考一卷 21 题）

在平面直角坐标系 $xOy$ 中，已知点
\[
F_1(-\sqrt{17},0),\qquad F_2(\sqrt{17},0),
\]
点 $M$ 满足
\[
|MF_1|-|MF_2|=2.
\]
记 $M$ 的轨迹为 $C$。

\begin{enumerate}
  \item 求 $C$ 的方程；
  \[
  x^2-\frac{y^2}{16}=1\qquad (x>0).
  \]

  \item 设点 $T$ 在直线
  \[
  x=\frac12
  \]
  上，过 $T$ 的两条直线分别交 $C$ 于 $A,B$ 两点和 $P,Q$ 两点，且
  \[
  |TA|\cdot|TB|=|TP|\cdot|TQ|.
  \]
  求直线 $AB$ 的斜率与直线 $PQ$ 的斜率之和。
\end{enumerate}

\end{document}
